\hypertarget{python-3.11-to-3.12-exception-groups-except-and-traceback-trees}{%
\chapter{Python 3.11 to 3.12: Exception Groups (except*) and Traceback
trees}\label{python-3.11-to-3.12-exception-groups-except-and-traceback-trees}}

\hypertarget{exception-groups-exceptiongroup-baseexceptiongroup-pep-654}{%
\subsection{\texorpdfstring{19.1 Exception Groups
(\texttt{ExceptionGroup} \& \texttt{BaseExceptionGroup} / PEP
654)}{19.1 Exception Groups (ExceptionGroup \& BaseExceptionGroup / PEP 654)}}\label{exception-groups-exceptiongroup-baseexceptiongroup-pep-654}}

Introduced in Python 3.11, \passthrough{\lstinline!ExceptionGroup!} and
\passthrough{\lstinline!BaseExceptionGroup!} allow raising and handling
multiple unrelated exceptions simultaneously. This is critical for
concurrent frameworks (such as \passthrough{\lstinline!asyncio!}
TaskGroups or Trio) where multiple background operations can crash
concurrently.

\hypertarget{exception-group-class-layout-and-hierarchy}{%
\subsubsection{1. Exception Group Class Layout and
Hierarchy}\label{exception-group-class-layout-and-hierarchy}}

The class hierarchy separates groups of standard exceptions from groups
containing system-level exceptions:

\begin{lstlisting}[language=Python]
class BaseExceptionGroup(BaseException):
    def __init__(self, message: str, exceptions: Sequence[BaseException]) -> None:
        self.message = message
        self.exceptions = list(exceptions)

class ExceptionGroup(BaseExceptionGroup, Exception):
    pass
\end{lstlisting}

\begin{itemize}
\tightlist
\item
  \passthrough{\lstinline!BaseExceptionGroup!} inherits directly from
  \passthrough{\lstinline!BaseException!}. It can wrap any exception
  instance, including system-terminating errors like
  \passthrough{\lstinline!KeyboardInterrupt!},
  \passthrough{\lstinline!SystemExit!}, or
  \passthrough{\lstinline!GeneratorExit!}.
\item
  \passthrough{\lstinline!ExceptionGroup!} inherits from both
  \passthrough{\lstinline!BaseExceptionGroup!} and
  \passthrough{\lstinline!Exception!}. It can only wrap exceptions that
  inherit from \passthrough{\lstinline!Exception!}. If a program
  attempts to instantiate an \passthrough{\lstinline!ExceptionGroup!}
  containing a \passthrough{\lstinline!BaseException!} that is not an
  \passthrough{\lstinline!Exception!} subclass, CPython raises a
  \passthrough{\lstinline!TypeError!} at runtime.
\end{itemize}

\hypertarget{cpython-c-struct-representation}{%
\subsubsection{2. CPython C-Struct
Representation}\label{cpython-c-struct-representation}}

At the C level, exception groups are managed by the
\passthrough{\lstinline!PyExceptionGroupObject!} structure, which
extends the standard exception header to support the nested tree
structure:

\begin{lstlisting}[language=C]
/* Include/cpython/exceptions.h */
typedef struct {
    PyBaseExceptionObject base; /* Base exception header containing dict, traceback, context */
    PyObject *msg;              /* Message describing the group (PyUnicodeObject) */
    PyObject *excs;             /* Tuple containing the child exception objects (PyTupleObject) */
} PyExceptionGroupObject;
\end{lstlisting}

When an exception group is instantiated: 1.
\passthrough{\lstinline!base.args!} is populated with a tuple containing
the \passthrough{\lstinline!message!} and
\passthrough{\lstinline!exceptions!} sequence. 2.
\passthrough{\lstinline!msg!} holds the string identifier directly for
fast attribute lookup. 3. \passthrough{\lstinline!excs!} holds the
underlying exceptions wrapped inside a flat, immutable tuple.

\begin{center}\rule{0.5\linewidth}{0.5pt}\end{center}

\hypertarget{the-except-clause-and-exception-tree-filtering}{%
\subsection{\texorpdfstring{19.2 The \texttt{except*} Clause and
Exception Tree
Filtering}{19.2 The except* Clause and Exception Tree Filtering}}\label{the-except-clause-and-exception-tree-filtering}}

To handle individual branches of an exception group, Python 3.11
introduced the \passthrough{\lstinline!except*!} statement. Unlike a
traditional \passthrough{\lstinline!except!} statement, which catches a
single exception instance, \passthrough{\lstinline!except*!} extracts a
matching subset from the Exception Group hierarchy, letting unmatched
exceptions propagate.

\hypertarget{compilation-and-bytecode-mechanics}{%
\subsubsection{1. Compilation and Bytecode
Mechanics}\label{compilation-and-bytecode-mechanics}}

The compiler generates a specialized sequence of bytecodes for
\passthrough{\lstinline!except*!} blocks to split the exception tree
dynamically at runtime.

Consider the two patterns below:

\begin{lstlisting}[language=Python]
# Standard try/except
try:
    raise ValueError("Error")
except ValueError as e:
    pass
\end{lstlisting}

Disassembly of standard \passthrough{\lstinline!except!}:

\begin{lstlisting}
  1           0 NOP
              2 SETUP_FINALLY           8 (to 12)

  2           4 LOAD_GLOBAL              1 (ValueError)
              6 LOAD_CONST               1 ('Error')
              8 CALL                     1
             10 RAISE_VARARGS            1

  3     >>   12 PUSH_EXC_INFO
             14 CHECK_EXCEPT             1 (ValueError)
             16 POP_JUMP_IF_FALSE       14 (to 46)
             18 STORE_FAST               0 (e)
             ...
\end{lstlisting}

Now, consider the \passthrough{\lstinline!except*!} variant:

\begin{lstlisting}[language=Python]
# Modern try/except*
try:
    raise ExceptionGroup("Main Group", [ValueError("Error A"), TypeError("Error B")])
except* ValueError as eg:
    print("Caught Val:", eg.exceptions)
\end{lstlisting}

Disassembly of \passthrough{\lstinline!except*!}:

\begin{lstlisting}
  1           0 NOP
              2 SETUP_FINALLY          12 (to 16)

  2           4 LOAD_GLOBAL              1 (ExceptionGroup)
              6 LOAD_CONST               1 ('Main Group')
              8 ... (building ValueError and TypeError objects)
             10 BUILD_LIST               2
             12 CALL                     2
             14 RAISE_VARARGS            1

  3     >>   16 PUSH_EXC_INFO
             18 CHECK_EXCEPT_STAR        1 (ValueError)
             20 POP_JUMP_IF_FALSE       16 (to 52)
             22 STORE_FAST               0 (eg)
             ...
\end{lstlisting}

\hypertarget{bytecode-analysis-and-vm-stack-transitions}{%
\subsubsection{2. Bytecode Analysis and VM Stack
Transitions}\label{bytecode-analysis-and-vm-stack-transitions}}

The crucial bytecode introduced for PEP 654 is
\passthrough{\lstinline!CHECK\_EXCEPT\_STAR!}. Its execution flow
proceeds as follows:

\begin{lstlisting}
Stack State Before CHECK_EXCEPT_STAR:
+------------------------+
|  ValueError (type tag)  | <-- TOP OF STACK (TOS)
+------------------------+
|  ExceptionGroup Object  | <-- TOS1
+------------------------+

Execution details of CHECK_EXCEPT_STAR:
1. Pops the target exception type (TOS: ValueError).
2. Inspects the active ExceptionGroup (TOS1).
3. Executes a C-level filtering function:
   - Splits the ExceptionGroup into a matched group containing all ValuerError instances.
   - Pushes the matched group.
   - Pushes the remaining unmatched exception group (containing TypeError).
   
Stack State After CHECK_EXCEPT_STAR (on match):
+------------------------+
| Matched ExceptionGroup | <-- TOS (Bound to 'eg' local variable)
+------------------------+
| Unmatched ExceptionGrp | <-- TOS1 (Propagated or checked by next handler)
+------------------------+
\end{lstlisting}

If the match group is empty (no \passthrough{\lstinline!ValueError!}
instances were found), \passthrough{\lstinline!CHECK\_EXCEPT\_STAR!}
pushes \passthrough{\lstinline!None!} to TOS, and
\passthrough{\lstinline!POP\_JUMP\_IF\_FALSE!} jumps directly to the
next handler block.

\begin{center}\rule{0.5\linewidth}{0.5pt}\end{center}

\hypertarget{exception-tree-filtering-algorithm}{%
\subsection{19.3 Exception Tree Filtering
Algorithm}\label{exception-tree-filtering-algorithm}}

The core mechanics of splitting exception groups are defined in
CPython's exception runtime library. The algorithm must recursively
inspect the tree of exceptions and cleanly separate them.

\begin{lstlisting}
                  ExceptionGroup ("Main")
                  /                     \
          TypeError("B")            ExceptionGroup ("Sub")
                                    /                     \
                             ValueError("A")       KeyError("C")
\end{lstlisting}

If we execute \passthrough{\lstinline!except* ValueError!}: 1. The
runtime calls the internal C function
\passthrough{\lstinline!exception\_group\_filter(eg, match\_value\_error\_func)!}.
2. It traverses the children of the group: *
\passthrough{\lstinline!TypeError("B")!}: Does not match. Added to the
\passthrough{\lstinline!unmatched!} collection. *
\passthrough{\lstinline!ExceptionGroup("Sub")!}: Recursively calls
\passthrough{\lstinline!exception\_group\_filter!}. * Inside
\passthrough{\lstinline!ExceptionGroup("Sub")!}: *
\passthrough{\lstinline!ValueError("A")!}: Matches! Added to the
sub-matched collection. * \passthrough{\lstinline!KeyError("C")!}: Does
not match. Added to the sub-unmatched collection. * Since there were
matches inside \passthrough{\lstinline!Sub!}, a new
\passthrough{\lstinline!ExceptionGroup("Sub")!} is instantiated,
containing only \passthrough{\lstinline!ValueError("A")!}. This is
returned as the sub-matched tree. * A new
\passthrough{\lstinline!ExceptionGroup("Sub")!} is instantiated
containing only \passthrough{\lstinline!KeyError("C")!} and returned as
the sub-unmatched tree. 3. The top-level execution collects the returned
structures: * \passthrough{\lstinline!matched!} tree:
\passthrough{\lstinline!ExceptionGroup("Main", [ExceptionGroup("Sub", [ValueError("A")])])!}
* \passthrough{\lstinline!unmatched!} tree:
\passthrough{\lstinline!ExceptionGroup("Main", [TypeError("B"), ExceptionGroup("Sub", [KeyError("C")])])!}

Here is the equivalent C-like pseudocode of the recursive filtering
function:

\begin{lstlisting}[language=C]
PyObject* exception_group_filter(PyObject *eg, PyObject *match_type) {
    PyObject *matched_list = PyList_New(0);
    PyObject *unmatched_list = PyList_New(0);
    
    PyObject *excs = ((PyExceptionGroupObject*)eg)->excs;
    Py_ssize_t size = PyTuple_GET_SIZE(excs);
    
    for (Py_ssize_t i = 0; i < size; i++) {
        PyObject *exc = PyTuple_GET_ITEM(excs, i);
        if (PyExceptionGroup_Check(exc)) {
            // Recursive split
            PyObject *sub_match = NULL, *sub_unmatch = NULL;
            split_exception_group(exc, match_type, &sub_match, &sub_unmatch);
            if (sub_match != NULL) {
                PyList_Append(matched_list, sub_match);
            }
            if (sub_unmatch != NULL) {
                PyList_Append(unmatched_list, sub_unmatch);
            }
        } else {
            // Leaf node matching
            if (PyErr_GivenExceptionMatches(exc, match_type)) {
                PyList_Append(matched_list, exc);
            } else {
                PyList_Append(unmatched_list, exc);
            }
        }
    }
    
    PyObject *matched_group = NULL;
    if (PyList_Size(matched_list) > 0) {
        matched_group = create_exception_group_from_list(eg, matched_list);
    }
    PyObject *unmatched_group = NULL;
    if (PyList_Size(unmatched_list) > 0) {
        unmatched_group = create_exception_group_from_list(eg, unmatched_list);
    }
    
    return PyTuple_Pack(2, matched_group, unmatched_group);
}
\end{lstlisting}

\begin{center}\rule{0.5\linewidth}{0.5pt}\end{center}

\hypertarget{traceback-representation-and-add_note-internals}{%
\subsection{\texorpdfstring{19.4 Traceback Representation and
\texttt{add\_note()}
Internals}{19.4 Traceback Representation and add\_note() Internals}}\label{traceback-representation-and-add_note-internals}}

\hypertarget{traceback-trees}{%
\subsubsection{1. Traceback Trees}\label{traceback-trees}}

Because exception groups represent hierarchical trees of errors,
CPython's traceback generator is modified to format tracebacks
recursively as tree diagrams.

When printed, the output represents the nesting level using visual
indicators:

\begin{lstlisting}
  + Exception Group Traceback (most recent call last):
  |   File "example.py", line 4, in <module>
  |     raise ExceptionGroup("Main Group", [ValueError("Error A"), TypeError("Error B")])
  | ExceptionGroup: Main Group (2 sub-exceptions)
  +-+---------------- 1 ----------------
    | ValueError: Error A
    +---------------- 2 ----------------
    | TypeError: Error B
    +-----------------------------------
\end{lstlisting}

If there are nested exception groups, the branch indicators prefix the
output recursively (e.g.,
\passthrough{\lstinline!+-+---------------- 1.1 ----------------!}).

\hypertarget{pep-678-exception-notes-and-c-level-internals}{%
\subsubsection{2. PEP 678 Exception Notes and C-Level
Internals}\label{pep-678-exception-notes-and-c-level-internals}}

PEP 678 introduces
\passthrough{\lstinline!BaseException.add\_note(note)!}, enabling users
to attach custom text to exceptions without modifying their
instantiation arguments. This is highly valuable for diagnostic tools,
test frameworks (like pytest), and tracing asynchronous executions.

When called, \passthrough{\lstinline!add\_note()!} stores the string
inside a \passthrough{\lstinline!\_\_notes\_\_!} list attribute on the
exception object. The C-level implementation handles memory allocations
and type checks directly:

\begin{lstlisting}[language=C]
/* Objects/exceptions.c */
static PyObject *
BaseException_add_note(PyBaseExceptionObject *self, PyObject *note)
{
    if (!PyUnicode_Check(note)) {
        PyErr_SetString(PyExc_TypeError, "note must be a string");
        return NULL;
    }

    PyObject *dict = self->dict;
    if (dict == NULL) {
        dict = PyDict_New();
        if (dict == NULL) {
            return NULL;
        }
        self->dict = dict;
    }

    PyObject *notes = PyDict_GetItemWithError(dict, &_Py_ID(__notes__));
    if (notes == NULL) {
        if (PyErr_Occurred()) {
            return NULL;
        }
        notes = PyList_New(0);
        if (notes == NULL) {
            return NULL;
        }
        if (PyDict_SetItem(dict, &_Py_ID(__notes__), notes) < 0) {
            Py_DECREF(notes);
            return NULL;
        }
        Py_DECREF(notes);
    }

    if (PyList_Append(notes, note) < 0) {
        return NULL;
    }

    Py_RETURN_NONE;
}
\end{lstlisting}

At runtime, the traceback printing logic calls
\passthrough{\lstinline!PyObject\_GetAttr(exc, \&\_Py\_ID(\_\_notes\_\_))!}.
If the list is found, it iterates over each note string and prints it
directly after the exception traceback and type representation.

\begin{center}\rule{0.5\linewidth}{0.5pt}\end{center}

\begin{center}\rule{0.5\linewidth}{0.5pt}\end{center}
